% GENERATED by build_tex.py from paper2/draft/paper2_draft.md (Appendix F) -- do not edit by hand

\section{Supplementary results}\label{app:F}

These passages were moved from the main text for length; they carry the same evidence and the same labels.

\subsection{Freezing the second layer (exploratory arm)}\label{app:F1}

Freezing the second layer instead of the first leaves a non-convex problem. The synthetic models still certified (16 of 16), but the real first layer grew by factors of $10^{4}$ to $10^{10}$, mostly through cancellation across its edges, and its endpoints were highly degenerate, with 21--41 near-flat directions; one endpoint that passed the certificate turned out, on a longer restart, to be still drifting outward (exploratory arm; post-hoc follow-up).

\subsection{Structured alternatives}\label{app:F2}

\textbf{Structured alternatives.} We tested two architectures designed to remove runaway routes by construction, against the early-stopped baseline, with leave-one-co-solvent-out (LOCO) validation:

\begin{table}[h]
  \centering
  % CAPTION-DRAFT: review
  \caption{Structured alternatives against the early-stopped baseline.}
  \label{tab:alternatives}
  \small
  \begin{tabular}{l p{0.24\linewidth} l l p{0.16\linewidth}}
  \toprule
  Model & Idea & LOCO RMSE & Converged & Seed agreement \\
  \midrule
  Baseline KAN-$\delta$ & As above & 0.255 & 0/10 & 0.112 \\
  Shared-curve KAN & One curve per input, mixed by penalised weights & 0.292 & 10/10 & 0.187 early; 0.085 polished \\
  Two-anchor excess model & Pure-solvent offset plus $x(1 - x)$ mixing term & 0.283 & not tested & not tested \\
  \bottomrule
  \end{tabular}
\end{table}

Seed agreement is the median prediction difference between seeds (lower is better). The shared-curve model converged every time, ending with no near-flat directions where the baseline had a median of 106 (range 86--116, at non-stationary points; Appendix~\ref{app:supp}, Figure~\ref{fig:uniqueness}a), and it was less accurate. At the early-stopped state its seeds agreed less than the baseline's (the pre-registered comparison); after polishing to certified minima their median disagreement fell to 0.085, with some seed pairs reaching the same minimum (post-hoc). The excess model's "pure co-solvent offset" tracked viscosity rather than dielectric constant (rank correlation $-$0.85, post hoc), reflecting that no training row contains a pure co-solvent. The shared-curve model's learned viscosity dependence also had the wrong sign on more than 10\% of rows in 3 of 10 seeds (a pre-registered prediction that failed). Chemically, conductivity should fall with viscosity at fixed composition, but in these data viscosity and dielectric constant rise together (r = 0.966), and higher permittivity favours salt dissociation; with the two nearly collinear, the data cannot separate their partial effects, which is the kind of non-identifiability this paper is about.

\subsection{The regularised teacher and gauge fixing}\label{app:F3}

\textbf{A known solution becomes recoverable.} Under the regularised loss, the original teacher is no longer a minimum. Optimising it under the same loss gives a "regularised teacher" computing nearly the same function (correlation 0.9992). Students match its units at 0.87--0.90 and its edges at 0.97--0.98, against 0.68--0.72 and 0.84--0.91 for the original teacher. This comparison was defined after the results and is exploratory.

\textbf{The penalty acts partly as gauge fixing.} At the certified minima, the layer-2 constant shift, an essentially exact symmetry, owes 96\% of its curvature to $\mu Q$: the regulariser selects one representative. The layer-1 constant shift (46\%) and unit rescaling (24\%) are only approximate symmetries in this architecture, because Gaussian bumps cannot represent constants exactly and tanh breaks rescaling.

\subsection{What a short path cannot tell}\label{app:F4}

\textbf{What a short path cannot tell.} We also predicted that a short regularisation path ($\gamma$ = $10^{-3}$, $10^{-5}$, $10^{-7}$), summarised by D = $\|\theta^{*}$($10^{-7}$)$\|^{2}$ / $\|\theta^{*}$($10^{-5}$)$\|^{2}$, would diverge more strongly as $\rho$ grows. It did not (Spearman 0.71 under P0 and 0.14 under P1, against the required 0.8): D was already high at every dose, including $\rho$ = 0. A post-hoc check, with its classification rule fixed before computing, explains why. Where the unregularised polish certified, mostly at low $\rho$, 30 of 32 paths were heading towards that certified finite minimum or its valley, none moved outward past it, and 2 found a better minimum at smaller norm. Along one direction with curvature $\lambda$, $D = \bigl((\lambda + 2\gamma_1)/(\lambda + 2\gamma_2)\bigr)^{2}$ with $\gamma_{1}$ = $10^{-5}$ and $\gamma_{2}$ = $10^{-7}$, so a finite minimum whose curvature lies within or near the window 2$\cdot10^{-7}$ to 2$\cdot10^{-5}$ keeps D large while the path is still approaching it; the observed low-$\rho$ medians (4.54 to 15.05) correspond to $\lambda$ = 6.7$\cdot10^{-6}$ to 1.7$\cdot10^{-5}$. A two-decade D therefore cannot separate a distant finite minimum from a receding one. Receding behaviour (Section 6) is shown only at $\rho$ = 0.95 and on the real data, over the tested range; at low $\rho$ the minima are finite but distant; at intermediate $\rho$ these tests do not settle the question. A secondary prediction, that the affine share of the parameter change would rise steadily with $\rho$, failed.

\subsection{The exact symmetry of B-spline edges}\label{app:F5}

Because B-splines sum exactly to one, constants can be moved between a unit's incoming edges and its bias without changing the loss: an exact 15-dimensional symmetry in the [4, 3, 1] model, which holds only approximately for RBF edges. It made the solver hang at converged points and would have made every certified point look degenerate. We excluded these analytically constructed directions, never fitted ones, from certification and from the degeneracy count; exact symmetries are redundancy in the parameterisation, while the runaway directions, which are approximate and data-dependent, remain subject to every check.

\subsection{The boundedness criterion}\label{app:F6}

Let L be continuous and bounded below, and R $\geq$ 0 continuous and coercive (R($\theta$) $\to$\allowbreak{} $\infty$ as $\|\theta\|$ $\to$\allowbreak{} $\infty$). Then $L + \gamma R$ is coercive for every $\gamma$ > 0 and has a global minimiser $\theta$\_$\gamma$. If L has a finite minimiser $\theta_{0}$, then L($\theta$\_$\gamma$) + $\gamma$R($\theta$\_$\gamma$) $\leq$ L($\theta_{0}$) + $\gamma$R($\theta_{0}$) and L($\theta_{0}$) $\leq$ L($\theta$\_$\gamma$), so R($\theta$\_$\gamma$) $\leq$ R($\theta_{0}$) for every $\gamma$, and coercivity keeps the $\theta$\_$\gamma$ bounded. Conversely, if the $\theta$\_$\gamma$ stay bounded as $\gamma \to 0$, they have a limit point $\theta$. For any $\theta$, L($\theta$\_$\gamma$) $\leq$ L($\theta$) + $\gamma$R($\theta$) $-$ $\gamma$R($\theta$\_$\gamma$) $\leq$ L($\theta$) + $\gamma$R($\theta$), so by continuity L($\theta$) $\leq$ L($\theta$): $\theta$ is a finite minimiser of L. The path in Section 6 follows certified local minima rather than global ones, so it provides evidence for the criterion's conclusion, not a proof of it.

\subsection{Robustness of the regularisation path (post hoc)}\label{app:F7}

Polishing from independent cold starts found no small-norm basin at any $\gamma$; at $\gamma$ = $10^{-7}$ under P1, 8 of 9 cold starts reached a \emph{lower} regularised objective than the warm path, at a larger norm. The warm path therefore understates the divergence. A straight line along the final drift direction rises almost at once: the minimisers recede along a curved valley, not along a ray. The growth sits in the penalty's blind directions, the affine RBF coefficients and SiLU weights of both layers, while the curved coefficients and the salt embedding shrink.

\subsection{Attribution, sign checks and edge agreement on fold A (post hoc)}\label{app:F8}

Three decompositions on the archived weights, with their plan committed before computing. (i) Splitting each polishing change into its blind-direction part (the edge-wise null-space projection of the RBF-coefficient change plus the SiLU-weight change) and the remainder, and applying each alone to the early-stopped model: the plan expected the blind part to be nearly free on training rows and costly on fold A. It was not: alone it raises both the ensembles' training and fold A RMSE above 2.9 (single seeds vary), and the remainder offsets most of that on training rows and less on fold A. On fold A the interaction is large and negative ($-$3.2 at $\mu$ = 0, $-$2.9 at $\mu$ = $10^{-6}$); on training rows it cannot be read as a share of the damage, because polishing lowers the training error. (ii) The sign of d log k/d(1000/T), row by row with the other inputs at their observed values: the anchor supplies the expected sign on almost every row, while polishing turns it on 25--33\% of fold A rows at the weaker penalties (0.3\% at early stopping, 5\% at $\mu$ = $10^{-4}$), concentrated at low temperature and high $x_{\mathrm{co}}$; the viscosity sign shows no clean ordering across states, although at $\mu$ = $10^{-4}$ its fold A violations sit mostly in the highest-$x_{\mathrm{co}}$ third (185 of 217 rows). (iii) Correlations of matched edge functions across seeds, with and without each edge's blind-direction part. The tables are generated from the results file.

\begin{table}[h]
  \centering
  % CAPTION-DRAFT: review (generated from paper2/revision/review4\_results.json)
  \caption{Post hoc: attribution of the fold A damage. RMSE (log k) of the 5-seed ensemble and of each seed, applying the blind-direction part or the remainder of the polishing change alone to the early-stopped model. On training rows the total change is negative, so no shares are given there.}
  \label{tab:r4attr}
  \footnotesize\setlength{\tabcolsep}{3pt}
  \begin{tabular}{l l r p{0.27\linewidth} r p{0.27\linewidth}}
  \toprule
   & State & Fold A & Fold A, seeds 0--4 & Training & Training, seeds 0--4 \\
  \midrule
  $\mu = 0$ & early & 0.1420 & 0.143, 0.137, 0.146, 0.199, 0.155 & 0.1290 & 0.134, 0.133, 0.141, 0.123, 0.129 \\
   & early + blind & 2.9110 & 1.053, 8.166, 0.760, 6.218, 1.869 & 3.6268 & 1.731, 11.573, 1.949, 8.481, 3.292 \\
   & early + remainder & 0.2141 & 1.304, 0.209, 0.203, 0.931, 1.044 & 0.2629 & 1.435, 0.298, 0.334, 0.847, 1.027 \\
   & polished & 0.8147 & 0.773, 0.968, 1.264, 1.337, 1.388 & 0.0801 & 0.084, 0.084, 0.072, 0.087, 0.087 \\
   & \emph{fold A shares} & blind 4.12 & remainder 0.11, interaction -3.22 &  &  \\
  $\mu = 10^{-6}$ & early & 0.1420 & 0.143, 0.137, 0.146, 0.199, 0.155 & 0.1290 & 0.134, 0.133, 0.141, 0.123, 0.129 \\
   & early + blind & 3.4938 & 0.471, 6.622, 6.300, 5.398, 1.746 & 3.0119 & 1.407, 7.487, 3.585, 6.741, 1.501 \\
   & early + remainder & 0.2344 & 0.980, 0.441, 0.550, 1.074, 0.751 & 0.3054 & 1.089, 0.501, 0.642, 1.264, 0.823 \\
   & polished & 1.0262 & 0.573, 0.739, 1.028, 0.883, 2.688 & 0.0845 & 0.086, 0.085, 0.085, 0.089, 0.087 \\
   & \emph{fold A shares} & blind 3.79 & remainder 0.10, interaction -2.90 &  &  \\
  \bottomrule
  \end{tabular}
\end{table}

\begin{table}[h]
  \centering
  % CAPTION-DRAFT: review (generated from paper2/revision/review4\_results.json)
  \caption{Post hoc: share of rows whose derivative has the unexpected (positive, $> 10^{-9}$) sign, 5-seed ensembles, overall and by thirds of $x_{\mathrm{co}}$ and of temperature (low to high) within each row set.}
  \label{tab:r4sign}
  \footnotesize\setlength{\tabcolsep}{3pt}
  \begin{tabular}{l l l r r r r r r r}
  \toprule
  Rows & State & Derivative & Share & $x_{\mathrm{co}}$ low & mid & high & $T$ low & mid & high \\
  \midrule
  training & early & $\partial\log k/\partial(1000/T)$ & 0.001 & 0.00 & 0.00 & 0.00 & 0.00 & 0.00 & 0.00 \\
   &  & $\partial\log k/\partial\ln\eta_{\mathrm{mix}}$ & 0.318 & 0.11 & 0.39 & 0.60 & 0.21 & 0.38 & 0.38 \\
   & $\mu = 0$ & $\partial\log k/\partial(1000/T)$ & 0.003 & 0.00 & 0.00 & 0.01 & 0.00 & 0.00 & 0.01 \\
   &  & $\partial\log k/\partial\ln\eta_{\mathrm{mix}}$ & 0.158 & 0.05 & 0.25 & 0.19 & 0.23 & 0.13 & 0.11 \\
   & $\mu = 10^{-6}$ & $\partial\log k/\partial(1000/T)$ & 0.003 & 0.00 & 0.00 & 0.01 & 0.00 & 0.00 & 0.01 \\
   &  & $\partial\log k/\partial\ln\eta_{\mathrm{mix}}$ & 0.265 & 0.18 & 0.35 & 0.28 & 0.32 & 0.25 & 0.22 \\
   & $\mu = 10^{-4}$ & $\partial\log k/\partial(1000/T)$ & 0.003 & 0.00 & 0.00 & 0.01 & 0.01 & 0.00 & 0.00 \\
   &  & $\partial\log k/\partial\ln\eta_{\mathrm{mix}}$ & 0.071 & 0.02 & 0.01 & 0.28 & 0.12 & 0.05 & 0.03 \\
  fold A & early & $\partial\log k/\partial(1000/T)$ & 0.003 & 0.00 & 0.00 & 0.01 & 0.01 & 0.00 & 0.00 \\
   &  & $\partial\log k/\partial\ln\eta_{\mathrm{mix}}$ & 0.143 & 0.00 & 0.00 & 0.49 & 0.18 & 0.14 & 0.10 \\
   & $\mu = 0$ & $\partial\log k/\partial(1000/T)$ & 0.249 & 0.17 & 0.22 & 0.39 & 0.59 & 0.09 & 0.00 \\
   &  & $\partial\log k/\partial\ln\eta_{\mathrm{mix}}$ & 0.072 & 0.01 & 0.05 & 0.19 & 0.19 & 0.00 & 0.01 \\
   & $\mu = 10^{-6}$ & $\partial\log k/\partial(1000/T)$ & 0.334 & 0.13 & 0.38 & 0.58 & 0.65 & 0.25 & 0.00 \\
   &  & $\partial\log k/\partial\ln\eta_{\mathrm{mix}}$ & 0.266 & 0.32 & 0.14 & 0.32 & 0.41 & 0.18 & 0.21 \\
   & $\mu = 10^{-4}$ & $\partial\log k/\partial(1000/T)$ & 0.054 & 0.00 & 0.03 & 0.15 & 0.15 & 0.00 & 0.00 \\
   &  & $\partial\log k/\partial\ln\eta_{\mathrm{mix}}$ & 0.342 & 0.00 & 0.17 & 0.99 & 0.36 & 0.34 & 0.33 \\
  \bottomrule
  \end{tabular}
\end{table}

\begin{table}[h]
  \centering
  % CAPTION-DRAFT: review (generated from paper2/revision/review4\_results.json)
  \caption{Post hoc: correlation of matched layer-1 edge functions across the 15 seed pairs (all 2069 training rows), the 8 continuous inputs and the combined salt edge, with and without each edge's blind-direction part; prediction-level $d_{AB}$ beside (certified: committed; early-stopped: computed with the committed definition).}
  \label{tab:r4edge}
  \footnotesize\setlength{\tabcolsep}{3pt}
  \begin{tabular}{l l r l r r r}
  \toprule
  Seeds & Version & Median & Range & $n$ & Salt edge & $d_{AB}$ \\
  \midrule
  early-stopped & full edges & +0.485 & -1.000 to +1.000 & 720 & +0.271 & 0.1220 \\
   & blind part removed & +0.013 & -1.000 to +1.000 & 720 & +0.142 &  \\
  certified $\mu = 10^{-4}$ & full edges & +0.372 & -1.000 to +1.000 & 720 & -0.190 & 0.0771 \\
   & blind part removed & +0.429 & -1.000 to +1.000 & 720 & -0.278 &  \\
  \bottomrule
  \end{tabular}
\end{table}
